\subsection*{5.2. 强 Lefschetz 定理 (LV)}

\subsubsection*{5.2.1.}
设 $S$ 为代数闭域 $k$ 上的 $n$ 维射影概形, $S\hookrightarrow\mathbb P$ 为射影嵌入. 以下式表示超平面截面的上同调类:
\[
u\in H^2(S,\mathbb Z_\ell(1))
\]
因而 $u$ 为 $\underline{\mathcal O}_{\mathbb P}(1)$ 在 $H^2(\mathbb P,\mathbb Z_\ell(1))$ 中的第一 Chern 类沿 $S\hookrightarrow\mathbb P$ 的拉回. 记
\[
L_S:H^\bullet(S,\mathbb Q_\ell(j))
   \longrightarrow H^{\bullet+2}(S,\mathbb Q_\ell(j+1))
\]
为与 $u$ 相乘的映射.

\subsubsection*{5.2.2.}
\underline{若对每个 (或某个) 整数 $j$ 和每个整数 $i>0$, 映射}
\[
L_S^i:H^{n-i}(S,\mathbb Q_\ell(j))
\longrightarrow H^{n+i}(S,\mathbb Q_\ell(i+j))
\]
\underline{都是同构, 则称 (配备 $u$ 的) $S$ 满足强 Lefschetz 定理 (LV).}

\subsubsection*{5.2.2.1.}
猜想当 $S$ 光滑时此条件成立. 当 $k$ 的特征为零时它成立 [6]. 因此, 当 $S$ 连同其射影嵌入能够提升到特征零时, 特化定理 [SGA 4, XVI 2.2] 给出同样的结论. 特别地, 它对光滑完全交成立. 对维数至多为二的光滑 $S$ [2], 以及对 Abel 簇, 也已知它成立; 对后者既可直接证明 (Lieberman [4] 和 [2]), 也可利用 Mumford 的定理 [5]: 极化 Abel 簇连同其极化可以提升到特征零.

回顾, 对射影簇 $S$ 及 $i\leq n$, 上同调的本原部分定义为
\begin{equation}\tag{5.2.2.2}
\operatorname{Prim}^i(S,\mathbb Q_\ell(j))
=\operatorname{Ker}\!\left(
L^{n-i+1}:H^i(S,\mathbb Q_\ell(j))
\longrightarrow
H^{2n-i+2}(S,\mathbb Q_\ell(j+n-i+1))
\right).
\end{equation}
若 $S$ 满足 (LV), 则当 $i\leq n$ 时有分解
\begin{equation}\tag{5.2.2.3}
H^i(S,\mathbb Q_\ell(j))
\xleftarrow[\sim]{\ 1+L\ }
\operatorname{Prim}^i(S,\mathbb Q_\ell(j))
\oplus H^{i-2}(S,\mathbb Q_\ell(j-1)).
\end{equation}

\subsection*{命题 5.2.3.}
\underline{设 $X$ 为光滑射影概形并满足 (LV). 设 $\tau:Y\hookrightarrow X$ 为一个光滑超平面截面的包含映射. 则}
\begin{equation}\tag{5.2.3.1}
\operatorname{Prim}^n(X,\mathbb Q_\ell(j))
=\operatorname{Ker}\!\left(
\tau^*:H^n(X,\mathbb Q_\ell(j))
\longrightarrow H^n(Y,\mathbb Q_\ell(j))
\right).
\end{equation}

\underline{证明.}
$\tau:Y\to X$ 的投影公式给出交换图
\begin{equation}\tag{5.2.3.2}
\begin{tikzcd}[column sep=huge,row sep=huge]
H^n(X,\mathbb Q_\ell(j))
  \arrow[rr,"L"]\arrow[dr,"\tau^*"']
  & {} &
H^{n+2}(X,\mathbb Q_\ell(j+1))
  \arrow[dl,"\tau_*"]\\
{} & H^n(Y,\mathbb Q_\ell(j)) & {}
\end{tikzcd}
\end{equation}
由对偶性与弱 Lefschetz 定理,
\[
\tau_*:H^n(Y,\mathbb Q_\ell(j))
\longrightarrow H^{n+2}(X,\mathbb Q_\ell(j+1))
\]
为同构的转置, 因此本身也是同构. 断言随之得出.

\subsection*{引理 5.2.4.}
\underline{设 $X$ 为光滑射影概形, $\tau:Y\hookrightarrow X$ 为一个光滑超平面截面的包含映射. 则 $X$ 满足 (LV) 当且仅当以下两个条件成立:}

\smallskip
\noindent
\textup{5.2.4.1.} \underline{$Y$ 满足 (LV).}

\smallskip
\noindent
\textup{5.2.4.2.} \underline{$H^{n-1}(Y,\mathbb Q_\ell)$ 上的杯积在子空间 $\operatorname{Im}H^{n-1}(X,\mathbb Q_\ell)$ 上的限制非退化.}

\underline{证明.}
有
\begin{equation}\tag{5.2.4.3}
\left\{
\begin{aligned}
L_X&=\tau_*\tau^*,\\[1ex]
L_Y&=\tau^*\tau_* \qquad\text{(由 4.2.11),}
\end{aligned}
\right.
\end{equation}
因而
\begin{equation}\tag{5.2.4.4}
L_X^{i+1}=\tau_*L_Y^i\tau^*.
\end{equation}
在交换图
\begin{equation}\tag{5.2.4.5}
\begin{tikzcd}[column sep=huge,row sep=huge]
H^{n-i-1}(X,\mathbb Q_\ell(j))
  \arrow[r,"L_X^{i+1}"]\arrow[d,"\tau^*"]
  &H^{n+i+1}(X,\mathbb Q_\ell(j+i+1))\\
H^{n-1-i}(Y,\mathbb Q_\ell(j))
  \arrow[r,"L_Y^i"]
  &H^{n-1+i}(Y,\mathbb Q_\ell(j+i))\arrow[u,"\tau_*"']
\end{tikzcd}
\end{equation}
中, 由弱 Lefschetz 定理, 当 $i\geq1$ 时 $\tau^*$ 为同构, 而由对偶性, $\tau_*$ 为同构. 因而 $X$ 满足 (LV) 当且仅当 $Y$ 满足 (LV) 且
\begin{equation}\tag{5.2.4.6}
L_X:H^{n-1}(X,\mathbb Q_\ell(j))
\longrightarrow H^{n+1}(X,\mathbb Q_\ell(j+1))
\end{equation}
为同构.

尚须将 (5.2.4.6) 与 5.2.4.2 对应起来. 设 $a,b\in H^{n-1}(X,\mathbb Q_\ell(j))$. 由 $X$ 上的 Poincaré 对偶性, (5.2.4.6) 等价于下列双线性型的非退化性:
\begin{equation}\tag{5.2.4.7}
\begin{aligned}
H^{n-1}(X,\mathbb Q_\ell(j))\times
H^{n-1}(X,\mathbb Q_\ell(j))
&\longrightarrow H^{2n}(X,\mathbb Q_\ell(2j+1)),\\
(a,b)&\longmapsto a\,L_X(b).
\end{aligned}
\end{equation}
但由 (5.2.4.3) 和投影公式,
\[
aL_X(b)=a\tau_*\tau^*(b)
=\tau_*\bigl(\tau^*(a)\tau^*(b)\bigr).
\]
这就证明了 (5.2.4.7) 和 5.2.4.2 的等价性.

\subsection*{推论 5.2.5.}
\underline{设 $X$ 为 $n$ 维光滑射影概形, $\tau:Y\hookrightarrow X$ 为一个光滑超平面截面的包含映射. 假设 $X$ 满足 (LV). 则当 $q\geq n$ 时, 映射}
\[
\tau^*:H^q(X,\mathbb Q_\ell(j))
\longrightarrow H^q(Y,\mathbb Q_\ell(j))
\]
\underline{为满射.}

\underline{证明.}
由 (5.2.4.3), 当 $i\geq0$ 时有
\[
L_Y^{i+1}=\tau^*L_X^i\tau_*,
\]
从而得到交换图
\begin{center}
\begin{tikzcd}[column sep=huge,row sep=huge]
H^{n-i}(X,\mathbb Q_\ell(j+1))
  \arrow[r,"L_X^i"]
  &H^{n+i}(X,\mathbb Q_\ell(j+i+1))\arrow[d,"\tau^*"]\\
H^{n-i-2}(Y,\mathbb Q_\ell(j))
  \arrow[u,"\tau_*"']\arrow[r,"L_Y^{i+1}"]
  &H^{n+i}(Y,\mathbb Q_\ell(j+i+1)).
\end{tikzcd}
\end{center}
由此得到所需的结果.

\subsection*{引理 5.2.6.}
\underline{设 $f:F\to B$ 为光滑固有态射, 且 $B$ 连通. 设 $\underline{\mathcal O}_F(1)$ 为 $F$ 上的相对丰沛可逆层, 并在每个几何纤维 $F_{\bar b}$ 上配备 $\underline{\mathcal O}_F(1)$ 所诱导的丰沛可逆层. 若某个几何纤维满足 (LV), 则所有其他几何纤维也满足 (LV).}

\underline{证明.}
设 $\xi_F\in H^2(F,\mathbb Z_\ell(1))$ 为 $\underline{\mathcal O}_F(1)$ 的第一 Chern 类. 此类及其幂 $\xi_F^i\in H^{2i}(F,\mathbb Z_\ell(i))$ 经由杯积诱导同态
\begin{equation}\tag{5.2.6.1}
\xi_F^i:R^{n-i}f_*\mathbb Q_\ell(j)
\longrightarrow R^{n+i}f_*\mathbb Q_\ell(j+i),
\end{equation}
其中 $n$ 为纤维的相对维数. 由固有基变换 [SGA 4, XII 5.2], 所有几何纤维满足 (LV) 当且仅当所有映射 (5.2.6.1) 都是同构. 由特化定理 [SGA 4, XVI 2.2], 层 $R^if_*\mathbb Q_\ell(j)$ 为扭常层. 因而映射 (5.2.6.1) 为同构当且仅当它们在某个给定几何点上为同构.

\subsection*{推论 5.2.7.}
\underline{设 $X$ 为光滑射影概形. 若某个光滑超平面截面满足 (LV), 则每个光滑超平面截面都满足 (LV).}

\underline{证明.}
$X$ 的光滑超平面截面的泛族满足引理 5.2.6 的假设.
